\section{1. Introduction}

\subsubsection{1.1 Research Background and Problem Significance}

Light field (LF) imaging has revolutionized computational photography by simultaneously capturing the spatial and angular distribution of light rays [?]. This rich 4D plenoptic function implicitly captures profound geometric and photometric cues, making LF depth estimation a cornerstone for downstream applications such as autonomous driving, augmented reality, and robotic manipulation [?]. Conventionally, the majority of LF depth estimation methods heavily rely on the Epipolar Plane Image (EPI) structure, exploiting the photo-consistency and linear disparity assumptions inherent in Lambertian scenes [?]. Under these ideal conditions, state-of-the-art algorithms have achieved remarkable success, demonstrating robust performance in purely diffuse environments.

However, obtaining the shape of glossy objects like metals, plastics, or ceramics remains challenging, since standard Lambertian cues like photo-consistency cannot be easily applied [?]. In real-world scenarios, LF images inevitably encounter non-Lambertian surfaces, including specular, transparent, and scattering materials. Inferring the depth of transparent or mirror (ToM) surfaces represents a hard challenge for either sensors, algorithms, or deep networks [?]. The fundamental difficulty lies in the fact that the physical interaction between light and non-Lambertian matter severely violates the view-consistency assumption. Specifically, the surface roughness, quantitatively defined as \[
r = a/\lambda
\] (where \(a\) is the surface height variation and \[
\lambda
\] is the wavelength), dictates the reflection type: \[
r \ll 1
\] leads to specular reflection, \[
r \approx 1
\] results in scattering, and \[
r \gg 1
\] yields diffuse reflection. Consequently, the Radiance Tensor Field (RTF) rank elevates from 1 in Lambertian scenes to \[
\geq 2
\] in non-Lambertian regions due to the superposition of multiple roughness levels and deflection vectors [?].

This physical violation leads to catastrophic performance degradation in existing methods. To illustrate the severity, our preliminary evaluations reveal that while a strong EPI-based baseline (e.g., EPINet) achieves an excellent Mean Absolute Error (MAE) of 0.081 in mixed Lambertian-dominant urban scenes, its MAE dramatically deteriorates to 0.411 in purely non-Lambertian scenes. Visual inspections explicitly confirm that in high-reflection or transparent regions, the network's predictions completely degenerate into mere superficial feature memorization of the input views, entirely failing to recover the underlying depth structures [1]. To mitigate this, some recent studies have attempted to incorporate defocus cues or employ data-driven Convolutional Neural Networks (CNNs) for material classification [?]. Unfortunately, defocus branches often fail to learn effective depth features in non-Lambertian regions, and purely data-driven 3-class CNN classifiers yield a suboptimal accuracy of merely 24.6\% on sparse \[
9 \times 9
\] angular patches, proving that limited patch information is inherently insufficient to support pure data-driven material parsing [?].

Therefore, there is an urgent and necessary demand to transcend the single Lambertian assumption and establish a unified physical framework capable of simultaneously handling diffuse, specular/scattering, and complex mixed scenes. On the one hand, the definition of depth itself might appear ambiguous in such cases; on the other hand, as humans can effortlessly deal with this through visual experience, depth sensing techniques based on physical optics hold the potential to address this challenge [?]. Rather than relying on heavily parameterized, black-box neural networks that struggle with sparse angular sampling, it is imperative to leverage the intrinsic angular frequency properties of the LF. By drawing an analogy to the k-space frequency analysis in Magnetic Resonance Imaging (MRI), where different tissues yield distinct frequency signatures, the 81 angular samples of an LF pixel can be transformed into the angular frequency domain to explicitly decode material-specific physical parameters [?]. This paradigm shift from data-driven heuristics to physics-driven angular frequency analysis is crucial for achieving robust, component-aware depth estimation in the wild.

\subsubsection{1.2 Limitations of Existing Methods}

Despite the substantial progress in light field (LF) depth estimation, inferring accurate geometry for non-Lambertian surfaces remains a formidable challenge. Existing approaches can be broadly categorized into traditional epipolar-based methods, early deep learning frameworks, and recent material-aware multi-cue networks. However, each category exhibits inherent limitations when confronting the complex physical interactions of light and matter, which can be systematically analyzed from both methodological and technical perspectives.

\#\#\#\# 1.2.1 Traditional Epipolar and Photo-Consistency Methods
Conventional LF depth estimation algorithms primarily exploit the Epipolar Plane Image (EPI) structure, relying heavily on photo-consistency and linear disparity assumptions [?]. While these methods achieve remarkable success in purely diffuse environments, they fundamentally struggle with non-Lambertian materials due to two critical flaws. 

First, the physical interaction between light and glossy, transparent, or scattering surfaces severely violates the view-consistency assumption. For instance, specular reflections and transparent refractions introduce view-dependent intensity variations, causing traditional matching costs to fail and resulting in severe erroneous texture replication artifacts or depth ambiguities [?]. Second, these methods lack an explicit physical model to characterize surface roughness and light deflection. They treat all pixels uniformly, failing to distinguish between diffuse, specular, and scattering components. Consequently, without quantifying the surface roughness \[
r = a/\lambda
\] or the wave vector variations, traditional methods cannot theoretically justify the failure of EPI structures in non-Lambertian regions, leaving the depth inference of such surfaces under-constrained.

\#\#\#\# 1.2.2 Early Deep Learning-based Approaches
To bypass the hand-crafted limitations of traditional methods, early deep learning-based approaches, such as EPINet and its variants, formulate depth estimation as an end-to-end regression task by implicitly learning EPI features [?]. Unfortunately, these data-driven models exhibit significant vulnerabilities in real-world applications. 

On the one hand, these networks demonstrate a strong bias toward the dominant Lambertian features present in training datasets. When deployed in mixed or purely non-Lambertian scenarios, they suffer from dramatic performance degradation due to poor generalization capabilities. Specifically, the networks tend to overfit to diffuse textures, leading to superficial feature memorization rather than genuine geometric reasoning in specular or transparent regions [?]. On the other hand, the purely data-driven nature of these models requires massive amounts of annotated non-Lambertian LF data, which is notoriously difficult and expensive to acquire, thereby limiting their scalability and robustness in unconstrained environments [?].

\#\#\#\# 1.2.3 Recent Material-Aware and Multi-Cue Networks
Recognizing the limitations of pure EPI and end-to-end regression, recent advancements have introduced material-aware and multi-cue frameworks. These methods attempt to integrate defocus cues, polarization information, or explicit material classification branches to assist depth estimation [?]. Despite their theoretical appeal, they encounter substantial bottlenecks in practice. 

Defocus-based branches, for example, often fail to extract reliable depth features in highly reflective or transparent regions where the point spread function (PSF) is heavily distorted by complex light transport [?]. Furthermore, methods relying on auxiliary material classification typically employ lightweight, data-driven CNNs on sparse angular patches (e.g., \[
9 \times 9
\]). As demonstrated in our preliminary analysis, such classifiers yield suboptimal accuracy (e.g., 24.6\%) because the limited spatial-angular receptive field is inherently insufficient to capture the high-order physical interactions governing non-Lambertian reflections [?]. Consequently, the erroneous material priors misguide the subsequent depth regression, compounding the estimation errors rather than mitigating them.

\subsubsection{1.3 Proposed Method Overview}

To overcome the aforementioned limitations, this paper proposes a novel physics-driven angular frequency analysis framework for robust LF depth estimation in complex mixed scenes. Instead of treating non-Lambertian reflections as noise or relying on fragile data-driven material classifiers, we explicitly model the light-matter interaction by transforming the sparse angular samples into the angular frequency domain. 

Drawing inspiration from k-space frequency analysis in MRI, where distinct tissues exhibit unique frequency signatures, we hypothesize that different material components (diffuse, specular, and scattering) possess distinguishable angular frequency distributions. By applying a 2D Fourier transform to the \[
9 \times 9
\] angular patch of each spatial pixel, we extract high-frequency and low-frequency angular features that intrinsically encode the surface roughness \[
r = a/\lambda
\] and light deflection vectors. These physics-derived frequency features are then integrated into a unified depth estimation network via a component-aware attention mechanism. This allows the network to dynamically adjust its receptive fields and matching costs based on the decoded physical parameters, ensuring accurate depth inference for both Lambertian and non-Lambertian regions without requiring explicit, error-prone material classification labels.

\subsubsection{1.4 Main Contributions}

The main contributions of this work are summarized as follows:

1. \textbf{Unified Physical Framework for Mixed Scenes}: We transcend the conventional single Lambertian assumption by establishing a unified physical framework that simultaneously accommodates diffuse, specular/scattering, and complex mixed scenes. This is achieved by explicitly incorporating the surface roughness and light deflection models into the depth estimation pipeline.
2. \textbf{Physics-Driven Angular Frequency Analysis}: We introduce a novel angular frequency domain analysis module, drawing an analogy to MRI k-space analysis. This module transforms sparse angular samples into the frequency domain to explicitly decode material-specific physical parameters, replacing unreliable data-driven heuristics with rigorous physical optics principles.
3. \textbf{State-of-the-Art Performance in Non-Lambertian Regions}: Extensive experiments on both synthetic and real-world LF datasets demonstrate that our proposed method significantly outperforms existing state-of-the-art algorithms. Specifically, it achieves remarkable improvements in depth accuracy for transparent, mirror, and glossy surfaces, effectively eliminating the erroneous texture replication artifacts prevalent in prior methods.